\chapter{从普通对称性到格拉斯曼代数}

\lectureoverview{超对称是一种对称性，因此我们首先回顾这个词在量子力学中的含义。物理变换在希尔伯特空间上由幺正算符表示；连续变换由无穷小生成，其交换子编码了变换如何复合。真正的对称性保持能量不变，这等价于与哈密顿量对易。超对称随后改变了所有熟悉例子共有的一个假设：其生成元把玻色子变为费米子，把费米子变为玻色子。构造这样的生成元自然导致产生与湮灭算符、反对易子，以及奇量平方为零的有限格拉斯曼代数。}

\section{对称性联系不同的位形}

对称性不仅仅是一种装饰性的规律性。它是一种把一个物理位形与另一个物理位形联系起来的操作。例如，左右反射把两只手联系起来。如果这个人以及周围的定律具有严格的左右对称性，那么两只手的对应性质就会一致。更一般地说，由严格对称性联系起来的物体或态必须共享该对称性所保持的物理量。在粒子物理学中，这常常意味着质量相等；在原子物理学中，这意味着处于同一转动多重态中的态具有相等的能量。

我们的目标是超对称，但先回顾转动和平移等连续对称性是有益的。超对称在通常意义上是否是“连续的”，这个问题本身就已经很棘手。它的无穷小参数不是普通的可交换实数。最好把它视为连续对称性的分级推广，而下面引入的格拉斯曼算术使这一陈述对我们目前的目的而言足够精确。\editorialnote{现代术语把超对称描述为一种连续分级对称性。这一限定解释了讲座中的警告：通常的连续与离散之分并不完全适用于超对称。}

\section{希尔伯特空间上的幺正操作}

量子系统由态矢量\(\ket{\psi}\)表示。物理对称操作由作用在该态上的算符\(U\)表示：
\begin{equation}
\ket{\psi}\longmapsto U\ket{\psi}.
\end{equation}
对于空间转动，转轴和转角决定\(U\)。如果一个态包含位于若干位置的粒子，那么变换后的右矢表示所有这些粒子都处于它们转动后的位置。平移的作用类似：它把该态替换为整个位形已被平移的态。

\begin{classroomqa}
\audiencequestion{态矢量属于一个常常是无穷维的希尔伯特空间。我们转动的是希尔伯特空间还是普通三维空间？\lecturetimestamp{00:04:20}}
\lecturerresponse{这两种描述在不同层面上指同一个操作。幺正变换作用在希尔伯特空间矢量上，并表示位形在普通空间中的物理转动。它并不是抽象希尔伯特空间字面意义上的三维转动；它是空间转动在希尔伯特空间中的表示。}
\end{classroomqa}

本质条件是幺正性：
\begin{equation}
U^\dagger U=UU^\dagger=\mathbf 1,
\qquad
U^{-1}=U^\dagger.
\end{equation}
这保持内积，从而保持概率：
\begin{equation}
\braket{U\psi}{U\psi}
=\bra{\psi}U^\dagger U\ket{\psi}
=\braket{\psi}{\psi}.
\end{equation}

\begin{figure}[H]
\centering
\includegraphics[width=0.56\textwidth]{lecture_04_figure_02.png}
\lectureframe{00:05:56.420}
\caption{黑板上从变换后的态到幺正性条件的早期过渡。右矢和带匕首的乘积可见，但讲课人遮住了它们的右端；完整的方程出现在下面的推导中。}
\end{figure}

严格地说，维格纳定理也允许反幺正对称性，例如时间反演。与恒等变换连续相连的变换，即本文的主题，由幺正算符表示。\editorialnote{反幺正的可能性是一项标准限定；它不改变讲座关于转动、平移或超对称生成元的讨论。}

\section{无穷小变换与生成元}

什么都不做由\(U=\mathbf 1\)表示，而不是\(U=0\)。足够小的变换与恒等变换相差一个关于其小参数线性的项。对于转过小角度\(\epsilon\)的转动，
\begin{equation}
U(\epsilon)=\mathbf 1+i\epsilon L+O(\epsilon^2).
\end{equation}
一阶幺正性要求
\begin{align}
U^\dagger U
&=\left(\mathbf 1-i\epsilon L^\dagger\right)
  \left(\mathbf 1+i\epsilon L\right)+O(\epsilon^2)\\
&=\mathbf 1+i\epsilon(L-L^\dagger)+O(\epsilon^2),
\end{align}
因而
\begin{equation}
L^\dagger=L.
\end{equation}
算符\(L\)就是无穷小生成元。

转动需要三个生成元，每个独立轴对应一个：
\begin{equation}
L_x,\qquad L_y,\qquad L_z.
\end{equation}
绕任意轴\(\widehat{\bm n}\)的转动由下式生成
\begin{equation}
L_{\widehat{\bm n}}=\widehat{\bm n}\cdot\bm L.
\end{equation}

\begin{classroomqa}
\audiencequestion{关于\(x\)、\(y\)和\(z\)轴的生成元能否相加，从而得到绕另一轴的转动？\lecturetimestamp{00:07:50}}
\lecturerresponse{可以。它们的线性组合生成绕中间轴的转动。三个独立分量足以构造出每一个无穷小空间转动。}
\end{classroomqa}

对于转动而言，这些生成元就是角动量的分量。这一记号在更一般的情形下也很有用，因为每一种连续对称性都可以通过其生成元在局部得到理解。

\section{作为乘法表的对易子}

转动生成元满足
\begin{align}
[L_x,L_y]&=iL_z,&
[L_y,L_z]&=iL_x,&
[L_z,L_x]&=iL_y,
\end{align}
或者，紧凑地写为，
\begin{equation}
[L_i,L_j]=i\epsilon_{ijk}L_k.
\end{equation}
当两个指标相同时，Levi-Civita符号为零；对于\(1,2,3\)的偶排列或奇排列，它分别为\(+1\)或\(-1\)。重要的结构性事实是封闭性：两个生成元的对易子仍是生成元。

为什么对易子能描述转动的几何？先绕\(x\)做一个小的转动，再绕\(y\)做一个小的转动，撤销\(x\)方向的转动，最后撤销\(y\)方向的转动。这些操作并不会相互抵消，因为绕不同轴的转动不对易。

\begin{classroomqa}
\audiencequestion{这种不对易性能否通过飞机的滚转、俯仰和偏航来想象？\lecturetimestamp{00:12:36}}
\lecturerresponse{可以。绕机体轴的相继转动提供了一个直接的几何演示，说明改变顺序会改变最终取向。下面的代数正是这一熟悉事实的无穷小版本。}
\end{classroomqa}

令
\begin{equation}
U_x(\epsilon)=\mathbf 1+i\epsilon L_x,
\qquad
U_y(\delta)=\mathbf 1+i\delta L_y.
\end{equation}
由于算符从右到左作用于右矢，所述的序列表示为
\begin{equation}
\mathcal C
=U_y^{-1}(\delta)U_x^{-1}(\epsilon)
 U_y(\delta)U_x(\epsilon).
\end{equation}

\begin{classroomqa}
\audiencequestion{为什么撤销操作可以写成逆而不是厄米共轭？\lecturetimestamp{00:13:25}}
\lecturerresponse{因为这些变换是幺正的：\(U^\dagger=U^{-1}\)。对于无穷小变换，取逆只需将一阶项的符号反转，\(U_x^{-1}=\mathbf 1-i\epsilon L_x+O(\epsilon^2)\)。}
\end{classroomqa}

在黑板上把四个因子相乘成了一次活跃的练习：所有关于\(\epsilon\)或\(\delta\)线性的项都抵消了，教室里的人追踪着混合项，而一次暂时的因子二混淆暴露出丢失算符顺序是多么容易。干净的二阶结果是
\begin{align}
\mathcal C
&=\mathbf 1+\epsilon\delta[L_x,L_y]
  +O(\epsilon^2\delta,\epsilon\delta^2)\\
&=\mathbf 1+i\epsilon\delta L_z
  +O(\epsilon^2\delta,\epsilon\delta^2).
\end{align}
因此，这个闭合序列留下了一个绕\(z\)的小转动，其角度正比于\(\epsilon\delta\)。\editorialnote{所展示的乘法在探索性的课堂计算之后确定了顺序和系数。反转序列的约定会反转符号，但不会改变残余变换由对易子支配这一结论。}

\begin{classroomqa}
\audiencequestion{改变逆操作的顺序会使乘积成为恒等变换吗？\lecturetimestamp{00:25:23}}
\lecturerresponse{如果每个变换都紧挨着它自己的逆放置，这些对就会抵消。非平凡的对易子序列有意将这些对分开；它未能坍缩这一点度量了顺序依赖性。}
\end{classroomqa}

更长的乘积确实会变得更复杂，但它们不包含任何新的原理。

\begin{classroomqa}
\audiencequestion{当生成元多于三个或有许多小变换时会发生什么？\lecturetimestamp{00:25:34}}
\lecturerresponse{会有更多的对易子和更多的记账工作，但知道了封闭的代数，就能确定任意多个无穷小变换如何组合。有限变换可以由许多小变换组装而成。}
\end{classroomqa}

在这个意义上，对易子代数就是连续群在恒等元附近写出的乘法表。

\section{平移补全了空间例子}

空间平移由动量生成。按照讲座中使用的符号约定，
\begin{equation}
\begin{aligned}
T_x(a)&=\mathbf 1+iaP_x+O(a^2),\\
T_y(b)&=\mathbf 1+ibP_y+O(b^2),
\end{aligned}
\end{equation}
其中\(a\)和\(b\)是距离而不是角度。

两个平移对易。将一点沿\(x\)方向平移，再沿\(y\)方向平移，撤销\(x\)方向的平移，再撤销\(y\)方向的平移；该点回到其原点：
\begin{equation}
[P_i,P_j]=0.
\end{equation}

\begin{classroomqa}
\audiencequestion{平移或转动一个物理对象是否需要力来启动和停止运动？\lecturetimestamp{00:28:45}}
\lecturerresponse{这并不是所定义的操作。对称变换是将原始状态与整个位形已被平移或转动的数学状态进行比较。它并不描述动力学轨迹、加速度或在移动装置时所做的功。}
\end{classroomqa}

旋转与平移不对易。平移一个点，再绕原点旋转它，然后沿原来的轴平移回去，并撤销旋转；这个点不会回到它出发的位置。混合代数关系为
\begin{equation}
[L_i,P_j]=i\epsilon_{ijk}P_k.
\end{equation}
这个关系表明，动量本身在旋转下按矢量变换。这里精确的指标练习不如一般性教训重要：欧几里得空间的变换被编码在其旋转与平移生成元之间的对易子中。

\section{对称性必须保持能量}

现在我们由变换的复合转向动力学。如果\(U\)确实是一个系统的对称性，那么施加它不能改变能量。设\(\ket{E}\)是一个能量本征态：
\begin{equation}
H\ket{E}=E\ket{E}.
\end{equation}
右边的符号\(E\)是一个数；\(\ket{E}\)是由该本征值标记的一个矢量。

\begin{figure}[H]
\centering
\begin{minipage}[t]{0.43\textwidth}
\centering
\includegraphics[width=\linewidth]{lecture_04_figure_03.png}
\lectureframe{00:34:38.239}
\end{minipage}
\hfill
\begin{minipage}[t]{0.50\textwidth}
\centering
\includegraphics[width=\linewidth]{lecture_04_figure_04.png}
\lectureframe{00:36:22.376}
\end{minipage}
\caption{变换后的右矢\(U\ket{E}\)与定义性的本征值方程\(H\ket{E}=E\ket{E}\)，分别隔离在连续的板书画面中。}
\end{figure}

如果\(U\)是一个对称性，那么\(U\ket{E}\)具有相同的本征值：
\begin{equation}
HU\ket{E}=EU\ket{E}.
\end{equation}
由于\(E\)是一个标量，它与\(U\)对易，并且可以使用原来的本征值方程：
\begin{align}
HU\ket{E}
&=EU\ket{E}\\
&=UE\ket{E}\\
&=UH\ket{E}.
\end{align}
因此
\begin{equation}
(HU-UH)\ket{E}=0.
\end{equation}

\begin{classroomqa}
\audiencequestion{这个推导能否从其第一个关系式重新进行？\lecturetimestamp{00:35:55}}
\lecturerresponse{从\(H\ket{E}=E\ket{E}\)开始。对称性表明\(U\ket{E}\)具有相同的能量，所以\(HU\ket{E}=EU\ket{E}\)。让标量\(E\)穿过\(U\)，将\(E\ket{E}\)替换为\(H\ket{E}\)，然后相减。结果是\((HU-UH)\ket{E}=0\)。}
\end{classroomqa}

\begin{figure}[H]
\centering
\includegraphics[width=0.70\textwidth]{lecture_04_figure_05.png}
\lectureframe{00:38:51.066}
\caption{完整的板书推导。弯曲的引导标记将等能量陈述连接到代换\(E\ket{E}=H\ket{E}\)；最后一行是\((HU-UH)\ket{E}=0\)。}
\end{figure}

这个论证对每一个能量本征态都成立。假设这些态构成一个完备基，那么湮灭其中每一个态的算符就是零算符：
\begin{equation}
[H,U]=0.
\end{equation}
这就是一个与时间无关的幺正变换成为哈密顿量的对称性的判据。

对于无穷小变换，
\begin{equation}
U(\epsilon)=\mathbf 1+i\epsilon L_i+O(\epsilon^2),
\end{equation}
恒等变换已经与\(H\)对易，所以每个生成元都必须满足
\begin{equation}
[L_i,H]=0.
\end{equation}

\begin{figure}[H]
\centering
\includegraphics[width=0.47\textwidth]{lecture_04_figure_06.png}
\lectureframe{00:42:14.296}
\caption{板上浓缩的对称性判据：\([U,H]=0\)，随后是其生成元形式，由部分模糊的下标标准化为\([L_i,H]=0\)。}
\end{figure}

如果\(L_i\)没有显式的时间依赖，海森堡方程给出
\begin{equation}
\frac{dL_i}{dt}
=i[H,L_i]=0
\qquad(\hbar=1).
\end{equation}
因此，对称性生成元是守恒荷。空间平移对称性给出动量守恒；旋转对称性给出角动量守恒。\editorialnote{显式时间导数的限定条件是标准做法，并且在讲座关于守恒律的陈述中是隐含的。}

\section{变换整个物理系统}

如果只变换装置的一部分，那么“旋转是一种对称性”这个说法就可能被误用。一个自由自旋可以被旋转而不改变其能量。然而，把它置于一个固定磁场中，并且只旋转自旋，它的塞曼能量就会改变。这并不否定旋转对称性；它意味着在自旋相对于场被旋转时，场被保持固定。把自旋和场一起旋转，或者使用一个本身不变的external field，能量就不变。

同样的问题出现在氢原子中。只让电子相对于质子平移会改变它们之间的间距，因而改变库仑能量。平移对称性作用于完整的电子-质子系统，使两者移动相同的量。

\begin{classroomqa}
\audiencequestion{外部场本身必须具有该对称性吗？\lecturetimestamp{00:43:09}}
\lecturerresponse{要么固定的外部场必须已经在该操作下不变，要么它必须与系统的其余部分一起被变换。自旋与未旋转的场之间的相对旋转是一种不同的物理构型，而不是整个装置的对称性变换。}
\end{classroomqa}

原子中的旋转简并给出了一个干净的正例。对于固定的角动量\(l\)，旋转混合磁量子数为
\begin{equation}
m=-l,-l+1,\ldots,l.
\end{equation}
的态。在没有外部磁场或电场时，多重态的所有成员都具有相同的能量：
\begin{equation}
E_{l,m}=E_l.
\end{equation}
对称性把不同的态联系起来，并迫使它们简并。

\section{超对称性改变了什么}

迄今为止讨论的每一种熟悉的对称性都保持玻色子与费米子之间的区别。旋转把一个电子变成另一个电子，而不是变成光子。色旋转混合夸克的三种颜色，但不会把夸克变成胶子。普通的内部对称性和时空对称性可以混合各分量或量子标签，然而它们的生成元是玻色的，不会改变费米子宇称。

超对称性则不同：
\begin{equation}
Q\ket{B}\sim\ket{F},
\qquad
Q\ket{F}\sim\ket{B}.
\end{equation}
它的多重态跨越了玻色子与费米子之间的分界。在一个严格的超对称理论中，同一多重态内的伙伴粒子具有相同的质量。这正是前几讲所要求的那种配对：匹配的玻色圈与费米圈可以抵消对标量质量的大紫外修正。这一构造难以直观想象，因为没有任何普通的空间运动能把整数自旋态变成半整数自旋态。因此我们采用代数方法来推进。

\section{一个原型超荷}

设\(A,A^\dagger\)湮灭和产生一个玻色子，并设\(C,C^\dagger\)湮灭和产生一个费米子。交换它们单粒子态的最简单的厄米算符是
\begin{equation}
Q=C^\dagger A+A^\dagger C.
\end{equation}
第一项移除一个玻色子并产生一个费米子；第二项移除一个费米子并产生一个玻色子。若\(Q\)作用在一个单玻色子态上，则\(A^\dagger C\)项为零，因为没有费米子可供湮灭，而\(C^\dagger A\)完成转换。作用在一个单费米子态上时，两者的角色互换。

\begin{classroomqa}
\audiencequestion{这一构造是否产生一个玻色子与费米子的叠加态这样的单个粒子？\lecturetimestamp{00:52:37}}
\lecturerresponse{\(Q\)是一个算符，而不是一个态。作用在一个纯玻色态上，它产生一个费米态；作用在一个纯费米态上，它产生一个玻色态。线性地作用在一个已经是一般叠加的态上，则产生另一个叠加态，正如任何线性算符所做的那样。}
\end{classroomqa}

下一间教室检验真空。由于两个湮灭算符都将其消灭，
\begin{equation}
A\ket{0}=0,\qquad C\ket{0}=0,
\end{equation}
我们有
\begin{equation}
Q\ket{0}=0.
\end{equation}
结果是零矢量，而不是真空的第二个副本。

这个原型是有意示意性的。一个自旋为零的玻色子可以映射到一个自旋为二分之一的费米子的任一分量，因此生成元本身需要一个旋量指标：
\begin{equation}
Q_\alpha\ket{B}\sim\ket{F_\alpha}.
\end{equation}
在讲座简化的两分量图像中，
\begin{equation}
Q_\alpha\ket{F_\beta}\sim\delta_{\alpha\beta}\ket{B}.
\end{equation}

\begin{classroomqa}
\audiencequestion{当一个自旋为零的玻色子有一个自旋分量而一个自旋为二分之一的费米子有两个时，态空间的维数如何能够匹配？\lecturetimestamp{00:55:02}}
\lecturerresponse{单个不带指标的\(Q\)是不够的。存在旋量分量\(Q_\alpha\)，每一个都能产生费米子的一个自旋分量。超荷本身带有半个单位的自旋，因为它使自旋改变半个单位。}
\end{classroomqa}

\(Q_\alpha\)是费米性的，或者更确切地说是格拉斯曼奇的算符。它们中的每一项都包含奇数个费米子算符，这就是它们反转费米子宇称的原因。普通的可观测物理量是偶的。应当把一个奇超荷与它的奇变换参数区分开来：一个有限的超对称变换由诸如\(\epsilon^\alpha Q_\alpha+\bar\epsilon_{\dot\alpha}\bar Q^{\dot\alpha}\)这样的偶组合构成，其中\(\epsilon\)和\(Q\)都是奇的。\editorialnote{讲座非正式地把生成元称为格拉斯曼数。更确切地说，超荷是奇算符，而反对易参数是格拉斯曼数；它们的乘积是偶的，可以出现在有限变换中。}

\section{玻色性与费米性算符代数}

对于以\(i,j\)标记的玻色模式，
\begin{align}
[A_i,A_j]&=0,&
[A_i^\dagger,A_j^\dagger]&=0,&
[A_i,A_j^\dagger]&=\delta_{ij}.
\end{align}
这些标记可以表示动量、位置、自旋或任何完备的单粒子基。由于产生算符对易，
\begin{equation}
A_i^\dagger A_j^\dagger\ket{0}
=A_j^\dagger A_i^\dagger\ket{0}.
\end{equation}
双玻色子态在交换下是对称的。

费米算符满足反对易关系：
\begin{align}
\{C_i,C_j\}&=0,&
\{C_i^\dagger,C_j^\dagger\}&=0,&
\{C_i,C_j^\dagger\}&=\delta_{ij},
\end{align}
其中
\begin{equation}
\{X,Y\}=XY+YX.
\end{equation}
对于单个模式，
\begin{equation}
\{C^\dagger,C^\dagger\}=2(C^\dagger)^2=0,
\end{equation}
所以
\begin{equation}
(C^\dagger)^2=0.
\end{equation}
该算符不为零；它的平方为零。它可以在该模式中产生一个费米子，但第二次作用给出零。这就是算符语言中的泡利不相容原理。

对于不同的标记，
\begin{equation}
C_i^\dagger C_j^\dagger\ket{0}
=-C_j^\dagger C_i^\dagger\ket{0},
\end{equation}
这就是双费米子态的反对称性。代数替换
\begin{equation}
\begin{gathered}
\text{bosons: commutators}\\
\Updownarrow\\
\text{fermions: anticommutators}
\end{gathered}
\end{equation}
编码了统计上的区别。

许多玻色子可以占据同一个模式。因此，一个大的光子占据数可以表现为具有普通数值振幅的经典电磁场。费米子不能在一个模式中累积，而一个奇的费米场不会获得普通的对易经典值。这促使我们为费米子的记账采用一种不同的算术。

\section{格拉斯曼数与宇称}

普通数可交换：
\begin{equation}
\alpha_i\alpha_j=\alpha_j\alpha_i.
\end{equation}
格拉斯曼生成元反交换：
\begin{equation}
\theta_i\theta_j=-\theta_j\theta_i.
\end{equation}
令\(i=j\)，在一个满足\(2\neq0\)的数系上，得到
\begin{equation}
\theta_i^2=0.
\end{equation}
不能给\(\theta_i\)赋予像六或七这样的普通数值。这些符号构成一个代数，其设计目的是追踪费米子符号。

普通系数与格拉斯曼生成元可交换：
\begin{equation}
\alpha\theta_i=\theta_i\alpha.
\end{equation}

\begin{classroomqa}
\audiencequestion{当一个普通数乘以一个格拉斯曼数时，得到的是何种对象？\lecturetimestamp{01:12:32}}
\lecturerresponse{它仍然是格拉斯曼奇的。乘以一个普通系数改变的是它的标度，而不是它的宇称。生成元的张成是普通数上的一个向量空间。}
\end{classroomqa}

该代数具有\(\mathbb Z_2\)宇称。普通系数以及偶数个生成元的乘积是偶的；一个生成元或奇数个生成元的乘积是奇的：
\begin{equation}
\begin{aligned}
\text{even}\cdot\text{even}&=\text{even},\\
\text{even}\cdot\text{odd}&=\text{odd},\\
\text{odd}\cdot\text{odd}&=\text{even}.
\end{aligned}
\end{equation}

\begin{classroomqa}
\audiencequestion{这是否类似于实数和虚数的乘法规则，还是像引入\(i\)那样的推广？\lecturetimestamp{01:14:25}}
\lecturerresponse{这是对算术的推广，但不应把这种类比推得太远。虚数仍然可交换并满足\(i^2=-1\)；格拉斯曼奇量反交换且平方为零。共同的奇偶乘法模式并不能使这些代数相同。}
\end{classroomqa}

像\(\theta_1\theta_2\)这样的乘积是偶的，并且与格拉斯曼代数的每个齐次元素可交换，但它不是普通数。它是非零的，而
\begin{equation}
(\theta_1\theta_2)^2=0.
\end{equation}

\begin{classroomqa}
\audiencequestion{如果\(\theta_1\theta_2\)是偶的，为什么它不只是一个普通数？\lecturetimestamp{01:17:36}}
\lecturerresponse{偶宇称意味着它在分次代数中可交换；这并不意味着它是实数或复数标量。它的幂零性立刻把它与一个非零普通数区分开来。}
\end{classroomqa}

\section{格拉斯曼变量的函数}

只有一个生成元时，每个多项式都在一阶之后终止：
\begin{equation}
f(\theta)=\alpha+\beta\theta.
\end{equation}
所有满足\(n\geq2\)的幂\(\theta^n\)都为零。有两个生成元时，最一般的元素是
\begin{equation}
f(\theta_1,\theta_2)
=\alpha+\beta_1\theta_1+\beta_2\theta_2
+\gamma\theta_1\theta_2.
\end{equation}
把最后一个单项式按相反顺序写出会改变其符号：
\begin{equation}
\theta_2\theta_1=-\theta_1\theta_2.
\end{equation}

\begin{classroomqa}
\audiencequestion{把\(\theta_1\)乘以一个系数\(\beta\)，会创造出第三个独立的格拉斯曼变量吗？\lecturetimestamp{01:18:39}}
\lecturerresponse{不会。它创造的是\(\theta_1\)的张成中的另一个向量，而不是一个新的基生成元。有两个独立生成元时，每个奇线性项都是\(\beta_1\theta_1+\beta_2\theta_2\)。}
\end{classroomqa}

对于三个生成元，有序基是
\begin{equation}
1,\quad
\theta_1,\theta_2,\theta_3,\quad
\theta_1\theta_2,\theta_1\theta_3,\theta_2\theta_3,\quad
\theta_1\theta_2\theta_3.
\end{equation}
一般来说，\(n\)个生成元产生\(2^n\)个有序基单项式。有限多个格拉斯曼生成元的任何解析函数都会截断为这样的多项式，因为重复的生成元为零。\editorialnote{\(2^n\)维基态明确展示了本讲对一、二、三个生成元所构造的模式。}

\section{格拉斯曼微分}

对于一个变量，定义
\begin{equation}
\frac{\partial}{\partial\theta}\alpha=0,
\qquad
\frac{\partial}{\partial\theta}\theta=1,
\end{equation}
从而
\begin{equation}
\frac{\partial}{\partial\theta}
(\alpha+\beta\theta)=\beta.
\end{equation}
对于两个变量，选择左导数并保持规范顺序\(\theta_1\theta_2\)：
\begin{align}
\frac{\partial^L f}{\partial\theta_1}
&=\beta_1+\gamma\theta_2,\\
\frac{\partial^L f}{\partial\theta_2}
&=\beta_2-\gamma\theta_1.
\end{align}
第二行中的负号出现，是因为奇导数必须先穿过\(\theta_1\)，然后才能对\(\theta_2\)求导。

这由分次乘积法则表达。如果\(A\)具有宇称\(|A|\in\{0,1\}\)，那么
\begin{equation}
\frac{\partial^L}{\partial\theta}(AB)
=\left(\frac{\partial^L A}{\partial\theta}\right)B
+(-1)^{|A|}A
\left(\frac{\partial^L B}{\partial\theta}\right).
\end{equation}
该规则在精神上是普通的，但包含反交换所要求的符号。格拉斯曼积分也可以定义，并且由于函数空间是有限的，它归结为一张短表。本讲推迟了那张表。

\section{指数谜题，已解决}

对于单个格拉斯曼生成元和普通系数\(a\)，
\begin{equation}
e^{a\theta}
=1+a\theta,
\end{equation}
因为所有更高次幂都为零。涉及同一生成元的两个指数满足
\begin{align}
e^{a\theta}e^{b\theta}
&=(1+a\theta)(1+b\theta)\\
&=1+(a+b)\theta\\
&=e^{(a+b)\theta}.
\end{align}

接着，教室里尝试不同的生成元。直接相乘给出
\begin{equation}
e^{\theta_1}e^{\theta_2}
=(1+\theta_1)(1+\theta_2)
=1+\theta_1+\theta_2+\theta_1\theta_2.
\end{equation}
起初，黑板上的展开暗示\(\theta_1+\theta_2\)的平方可能提供最后一项。一位听众抓住了符号：
\begin{equation}
(\theta_1+\theta_2)^2
=\theta_1^2+\theta_1\theta_2
 +\theta_2\theta_1+\theta_2^2
=0.
\end{equation}

\begin{classroomqa}
\audiencequestion{平方项是否因为它包含\(\theta_1\theta_2+\theta_2\theta_1\)而抵消？\lecturetimestamp{01:24:00}}
\lecturerresponse{是的。这一修正改变了结论：\(e^{\theta_1+\theta_2}=1+\theta_1+\theta_2\)，而有序乘积\(e^{\theta_1}e^{\theta_2}\)也包含\(\theta_1\theta_2\)。朴素的指数加法法则不适用于不同的反交换指数。}
\end{classroomqa}

这并不矛盾。即便是普通矩阵，也只有在\(A\)与\(B\)对易时，才满足\(e^Ae^B=e^{A+B}\)。这里
\begin{equation}
[\theta_1,\theta_2]
=\theta_1\theta_2-\theta_2\theta_1
=2\theta_1\theta_2.
\end{equation}
这个对易子是偶的，并且与两个生成元都对易，因此Baker-Campbell-Hausdorff级数终止：
\begin{equation}
e^{\theta_1}e^{\theta_2}
=\exp\left(
\theta_1+\theta_2+\frac12[\theta_1,\theta_2]
\right)
=e^{\theta_1+\theta_2+\theta_1\theta_2}.
\end{equation}
由于右端的指数是幂零的，它恰好展开为
\(1+\theta_1+\theta_2+\theta_1\theta_2\)。\editorialnote{这个BCH计算给出了现场讲座留作问题的解答。讲座正确地发现了朴素定律的失效，但在指出缺失的对易子项之前就结束了。}

\begin{classroomqa}
\audiencequestion{\(1+\theta_1\)是什么样的数，普通部分与Grassmann部分允许相加吗？\lecturetimestamp{01:26:16}}
\lecturerresponse{它是Grassmann代数中的一个非齐次元素，具有偶部分\(1\)和奇部分\(\theta_1\)。这样的和是允许的，但它们既不是普通标量，也不是齐次奇元素。}
\end{classroomqa}

接近末尾时，另一个表面上的悖论受到检验：
\begin{equation}
(\theta_1+\theta_2)^2=0.
\end{equation}
这并不意味着\(\theta_1\theta_2=0\)，因为两个交叉项是相互抵消，而不是各自消失。

\begin{classroomqa}
\audiencequestion{如果每个奇Grassmann元素的平方都为零，那是否意味着任意两个Grassmann生成元的乘积都为零？\lecturetimestamp{01:32:19}}
\lecturerresponse{不是。方程\((\theta_1+\theta_2)^2=0\)说的是\(\theta_1\theta_2+\theta_2\theta_1=0\)。它证明的是反对易性，而不是两个有序乘积中任何一个为零。不同的生成元可以有非零乘积。}
\end{classroomqa}

\begin{classroomqa}
\audiencequestion{当我们写下这些指数时，是否假设了普通的幂级数定义？\lecturetimestamp{01:34:42}}
\lecturerresponse{是的，而且这个定义是自洽的，因为每个级数都会终止。失效的不是指数本身，而是对非对易的\(A\)和\(B\)不成立的身份\(e^Ae^B=e^{A+B}\)。对易子修正恢复了正确的合成律。}
\end{classroomqa}

这个探索性的结尾是恰当的。普通对称性把我们由酉变换引向生成元和对易子；超对称性把我们由奇生成元引向反对易变量。在两种情形中，乘法律都是核心对象。对于转动，它出现在\([L_i,L_j]\)中；对于Grassmann指数，它出现在阻止朴素相加的对易子项中。这套代数并不熟悉，但它是自洽的、有限的，并且恰好适合超对称性所需的费米子变换。
